% Ad Familiares 16.12 --- headnote
% ad-familiares-16-12, 29 January 49 BC, Capua

\textit{Cicero to his freedman Tiro, written from Capua
on the fourth day before the Kalends of February 49 BC
(the manuscript dateline: \textit{Scr. Capuae iv K.
Febr. a. 705 (49)}). Tiro is still at Patrae,
recovering; the previous letter to him from Italy
(\textit{Fam.}~16.11) had reached him three weeks ago
with the first news of the crisis. This is the long
follow-up, written after Cicero has taken up the
Capuan command --- the assignment briefly named in
that earlier letter --- with the situation fully out
in the open.}

\textit{The shape of the letter is a full political
brief framed inside the affectionate Tiro-letter
opening and close. Sections 1--2 are the
self-justification: the city has been abandoned, and
Cicero, who from the moment he reached the walls had
been saying that nothing is worse than civil war, was
not heard. Section 3 lays out the terms now being
floated from Caesar's side --- Pompey to Spain,
levies disbanded, the Gauls handed over, Caesar to
seek the consulship in person; Section 4, the
alternative if the terms collapse, and the great
piece of news, that Labienus has come over. Section 5
is Cicero's own position --- the coastal command
beginning from Formiae, deliberately kept small so
that his pleas for peace can still carry weight --- and
the private grief that Dolabella is with Caesar. Section
6 hands Tiro over to A. Varro for the crossing and
returns to the medical-affectionate register. The
last words --- where Marcus is, where the women are
--- are the orientation Tiro will want first when he
opens it.}
